\chapter{Cum e construit, verificat și publicat}
\label{ch:tehnic}

Un ultim capitol, mai scurt, pentru cine vrea să știe și partea de
inginerie din spatele aplicației, nu e conținut de curs, dar explică de ce
demo-urile din capitolele anterioare pot fi luate ca atare.

\section{Pe scurt, arhitectura}

Front-end-ul e React peste Vite. Back-end-ul e FastAPI cu NumPy, SciPy și
PyWavelets. Fiecare ecran e o intrare într-un singur tablou de date, nu un
component separat, iar douăzeci și trei din cele cincizeci montează integral
o componentă de demo interactiv, aceleași componente pe care bara laterală
le poate deschide și individual, în afara turului. Toate desenează pe o
suprafață fixă de $1440\times 810$ pixeli virtuali, scalată proporțional la
fereastra reală, ca un ecran gândit pentru un videoproiector să nu piardă
nimic pe un monitor mai mic.

Calculul e împărțit în două, și distincția merită spusă limpede. Ecranele
care raportează cifre măsurate trimit parametrii către back-end și primesc
înapoi rezultatul unei biblioteci consacrate: filtrele, kernelele pe imagini
reale, denoisingul, comparația la bitrate egal, electrocardiograma și
extragerea benzilor EEG trec toate prin NumPy, SciPy și PyWavelets. Ecranele
care ilustrează un mecanism calculează în browser, cu filtre Haar scrise
direct în JavaScript: descompunerea pas cu pas, banca de filtre, piramida,
reconstrucția, scalograma, cutiile Heisenberg și etapele intermediare ale
pipeline-ului JPEG. Ele arată forma unei idei pe un semnal ales anume, nu o
măsurătoare pe date reale, și de aceea nu poartă cifre pe care textul să se
sprijine.

\section{Verificare automată}
\label{sec:gate}

Corectitudinea vizuală a celor cincizeci de ecrane se verifică prin trei
unelte separate, fiecare acoperind o clasă diferită de defect, în loc de o
singură trecere manuală.

Prima parcurge tot turul într-un browser real și verifică opt clase de
problemă pe fiecare ecran (Tabelul~\ref{tab:gate}), plus erorile de consolă
JavaScript. Conform ultimei rulări complete, verificarea e curată pe toate
cele cincizeci de ecrane, la două rezoluții diferite.

\begin{table}[htbp]
\centering
\caption{Cele opt clase de defect verificate automat pe fiecare ecran.}
\label{tab:gate}
\small
\begin{tabular}{@{}lp{10.3cm}@{}}
\toprule
Clasă & Ce detectează \\
\midrule
\txt{pageScroll} & Suprafața de design scrolează, conținutul depășește $1440\times 810$. \\
\txt{outOfFrame} & Un element vizibil (buton, input, imagine, canvas, text) iese din
  viewport, cu toleranță de 2 pixeli. \\
\txt{innerScroll} & Un element are conținut tăiat sau scrolabil intern
  (\txt{scrollHeight} peste \txt{clientHeight}). \\
\txt{overlaps} & Două elemente vizibile fără relație de descendență se
  suprapun cu cel puțin 6 pixeli pe ambele axe. \\
\txt{occludedControls} & Un control al cărui centru, testat la trei
  înălțimi, indică sub cursor un element diferit, prinde acoperirea de
  către un element exclus din \txt{overlaps}, de exemplu un \txt{canvas}. \\
\txt{smallMath} & O formulă KaTeX randată sub pragul de lizibilitate (14px
  inline, 18px afișată, în pixeli de design). \\
\txt{rawTokens} & Un token \texttt{\$...\$} sau \txt{[cite:...]} ajuns pe ecran ca
  text brut. \\
\txt{duplicateTitles} & Un ecran de tip \txt{embed} își afișează propriul
  titlu, deasupra titlului din metadata turului. \\
\bottomrule
\end{tabular}
\end{table}

A doua unealtă merge mai departe decât o simplă citire: apasă efectiv
fiecare buton, mișcă fiecare cursor, comută fiecare selector și bifează
fiecare căsuță de pe fiecare ecran, cu o captură înainte și una după fiecare
acțiune (Figura~\ref{fig:kernels-edu-interaction}). Măsoară dacă vreun
element s-a deplasat neașteptat pe ecran, verifică geometric dacă imaginile
și canvas-urile rămân în interiorul cutiei lor și colectează orice eroare
nouă de consolă. Ultima rulare completă a apăsat sau tras 911 controale, pe
toate cele cincizeci de ecrane, fără niciunul lipsit de un nume accesibil și
fără nicio eroare.

\begin{figure}[htbp]
\centering
\includegraphics[width=\textwidth]{15-kernels-edu.png}
\caption{Un ecran tipic pentru trecerea de interacțiune: selector de sprite,
  selector de kernel, cursor de dimensiune, selector de margine, cursor de
  viteză și cinci butoane de transport, fiecare apăsat sau tras automat, cu
  o captură înainte și una după.}
\label{fig:kernels-edu-interaction}
\end{figure}

A treia unealtă scanează codul sursă, nu ecranul: caută orice caracter
Unicode folosit în locul unei formule LaTeX sau al unei iconițe proprii de
interfață, pe principiul că matematica se randează mereu prin LaTeX, iar
orice iconiță e un desen propriu, nu un caracter de font care arată diferit
de la un sistem la altul. E curată pe tot front-end-ul.

Niciuna dintre cele trei unelte nu vede ce se desenează în interiorul unui
element de tip canvas, și acolo trăiesc majoritatea graficelor din tur:
spectre, răspunsuri de filtru, forme de wavelet, scalograme. O suprapunere
între o etichetă și o linie de grilă, acolo, nu apare în niciun raport
automat. Compensarea e directă: citirea capturilor de ecran salvate de
aceste unelte, ecran cu ecran, exact așa au fost alese și cele douăzeci și
două de figuri din acest document.

\section{Rulare locală}

Back-end-ul pornește din \txt{backend/}, după instalarea dependențelor
(FastAPI, Uvicorn, NumPy, SciPy, Pillow, PyWavelets, Pydantic). Front-end-ul
pornește din \txt{frontend/}, după \texttt{npm install}. Serverul de dezvoltare
Vite redirecționează automat orice cerere către \txt{/api} spre back-end,
așa că front-end-ul n-are niciun URL scris în cod, nici măcar în dezvoltare.
Fiecare ecran al turului se deschide direct, cu propriul \emph{deep link}.

\section{Publicarea pe server}

Aplicația rulează dintr-o singură imagine Docker, construită în doi pași:
front-end-ul se compilează întâi, apoi rezultatul e copiat peste back-end,
care ajunge astfel să servească, pe același port, și pagina, și API-ul.
Local, imaginea se construiește și pornește cu o singură comandă
\texttt{docker compose}.

Publicarea la \url{https://ale0204.student-dev.ro/dsp/} nu construiește nimic
pe mașina de acasă: sursa proiectului se împachetează și se trimite pe
server, imaginea se construiește direct acolo, deci mașina de acasă n-are
nevoie de Docker instalat sau pornit, doar trimite codul, iar la final
containerul pornește în spatele reverse proxy-ului serverului, care termină
conexiunea securizată și trimite mai departe cererile către aplicație.
Containerul nu publică niciun port propriu.

Fișierele care descriu un server anume, rețeaua în care intră containerul,
contul și directorul de pe gazdă, stau în afara depozitului public, pentru
că depind de gazdă, nu de aplicație. Depozitul rămâne astfel rulabil de
oricine, oriunde, fără nicio adresă scrisă în cod.
